\section{Los FPGA}
    \subsection{Definición}
        Son circuitos integrados diseñados con la arquitectura basada en CLBs y LUTs, programables mediante SRAM. Las siglas provienen de \textit{Field-Programmable Gate Array}, que por separado hacen referencia a:
        \begin{itemize}
            \item \textit{Field-Programmable}: El usuario puede reprogramar el integrado sin la necesidad de mandarlo de vuelta a la fábrica para aplicar cambios.
            \item \textit{Gate Array}: matriz de compuertas.
        \end{itemize}
        \sangria{} La propiedad de que sean reprogramables viene del primer FPGA fabricado por Xilinx cuyo diseño implementaba la SRAM (Static-RAM): una memoria que contiene las conexiones entre transistores. Antiguamente se usaba unos ''fusibles'', por lo que la reprogramación era imposible de llevar a cabo. Los términos CLB y LUT vienen de la SRAM:
        \begin{itemize}
            \item \textbf{LUT}: de \textit{Look-Up Table}, es una pequeña parte de la memoria SRAM con $2^N$ entradas que corresponden luego a las celdas de memoria, permitiendo representar cualquier función booleana porque cubre todas las combinaciones posibles según la cantidad de entradas que se den.\cite{harris}
            \item \textbf{CLB}: de \textit{Configurable Logic Block}, contiene uno o varios LUTs más un Flip-Flop. Los LUTs hacen la lógica combinacional, y el Flip-Flop agrega la capacidad de memoria/secuencial.
        \end{itemize}
    \subsection{Fabricantes destacados y software}
        \sangria{} Fue la empresa Xilinx quien inventó el primer FPGA utilizando la SRAM, y con comunicación con el exterior mediante GPIO puertos (\textit{General Porpose Input/Output}). Se programan en HDL, siendo ISE el usado por Xilinx hasta ser reemplazado por Vivado.
